% Section 6.3, from lectures/L06/appendix/b_exercises.md.

\section{Exercises}
\label{c6:sec:exercises}\label{c6:app:b}

\begin{aside}
\textbf{How to work these.} The Reflection and Calculation exercises are answered on paper, from
the chapter. The Code exercises are programs for the board, except Exercise~\ref{c6:ex:3}, which
is a plain PC program: build and flash the board programs the way \secref{c2:app:c} (Microchip
Studio) or \secref{c2:app:d} (Linux) describes, or run them in the simulator of
Appendix~\ref{sim:ch} if you have no board yet.

\textbf{How to check your work.} The Reflection and Calculation exercises are answered in
Appendix~\ref{app:answers}. The Code exercises have no printed solution: the board, or for
Exercise~\ref{c6:ex:3} the terminal, tells you whether a program does what the exercise describes.
Try each exercise before you read its answer.
\end{aside}

% ------------------------------------------------------------------------------------------
\exerciseset{Concepts}
\label{c6:set:1}

\calcexercise{\code{UBRR0} for 19200 baud}
\label{c6:ex:1}
Calculate the value of \code{UBRR0} for a baud rate of \code{19200} at a 16\,MHz clock. Round to
the nearest integer.

\reflexercise{What \code{UDRE0} means}
\label{c6:ex:2}
Explain what the \code{UDRE0} flag represents and why \code{serial_write_byte()} waits for it
before writing to \code{UDR0}.

% ------------------------------------------------------------------------------------------
\exerciseset{Program}
\label{c6:set:2}

\codeexercise{Decimal digits by hand}
\label{c6:ex:3}
Write a \code{serial_write_uint(const uint8_t byte)} helper that prints \code{byte} as decimal
text, one digit at a time, e.g.\ \code{serial_write_uint(243)} prints \code{243}. Extract each
digit yourself via division/modulo; don't just hand the whole number to
\code{printf("\%u", byte)}/\code{sprintf()} and let it do the conversion for you, that's the same
idea you'll need on any future toolchain where \code{sprintf()}/float formatting isn't available.

This is a console exercise: a plain PC program is enough, no board or UART required.
\code{printf()} (or \code{putchar()}) stands in for \code{serial_write_byte()}, one call per digit.

Hint: dividing/modulo-ing a number by \code{10} repeatedly gives you its digits least-significant
first; printing them in that order comes out backwards. Instead, find the largest power of ten
that still fits inside \code{byte} first (\code{uint8_t} maxes out at \code{255}, so at most 3
digits), then walk that divisor down by a factor of \code{10} each step, digits come out in the
right order with no need to reverse anything afterward. Watch out for \code{0}: make sure the
search for the starting divisor doesn't shrink it all the way down to nothing before the loop that
actually emits digits ever runs.

\codeexercise{Receiving a command}
\label{c6:ex:4}
Extend \file{serial_driver_demo}'s driver with receive support, following \secref{c6:a:7}'s
\code{RXEN0}/\code{RXC0} pattern: enable the receiver in \code{serial_init()}, then add a
\code{serial_read()} to \file{serial.h}/\file{serial.c} that blocks until a byte arrives.

Use it in \code{main()} so that receiving \code{'t'} over serial toggles LED1 the same way pressing
the button already does, reporting the new state the same way \code{main()} already reports a
button press.

\codeexercise{A temperature reporter}
\label{c6:ex:5}
Capstone, tying together ADC (Chapter~\ref{c4:ch}), timers (Chapter~\ref{c5:ch}), and UART (this
chapter). Read ambient temperature from a TMP36 sensor on A0 and report it over UART, either once a
minute (Timer-driven) or immediately on a button press on D13, debounced with a timer so contact
bounce doesn't spam the report:
\begin{itemize}
  \item Reuse \code{adc_read()} from \secref{c4:a:3} to sample A0.
  \item TMP36 outputs \code{750 mV} at \code{25 °C} and rises \code{10 mV} per \code{°C};
    converting the ADC's \code{0-1023} reading to a voltage and then to a temperature gives:

\begin{console}
T (°C) = 500 * adc_reading / 1023 - 50
\end{console}

  \item The reported message should look like \code{"Temperature: 23 C"} (build it with
    \code{sprintf()}, as in \secref{c6:a:5}).
\end{itemize}
For a harder version: keep the five most recent readings in a small array and report the running
average alongside the instantaneous reading; and instead of a fixed one-minute report interval,
track the time between the last five button presses and auto-adjust the report interval to match
how often the button is actually being pressed.
